\subsection{Global modularity tests for a same-discriminant unit}
\label{sec:hadamard-modularity-scope}

This is a bounded method audit, not a simultaneous-unit theorem.  Let
$p\ge31$ be prime with $p\equiv3\pmod4$ and put $b=(p+3)/2$.
The already proved full-series congruence
$H_b\equiv U_pH_2\pmod p$ turns the remaining coefficient target into
\[
 H(2,pd)H(3,d)\ne0\pmod p,
 \qquad -d<0\text{ odd fundamental},\quad d<p.
 \tag{HM.1}
\]
The first-unit theorem supplies such a $d$ for the first factor, and its
class number is already a unit.  It does not supply the second factor.
The operations below do not complete this step.

\paragraph{An exact same-index Dirichlet series is not yet an integral
modular detector.}
Mizuno's \emph{Rankin--Selberg convolutions of noncuspidal half integral
weight Maass forms in the plus space}, Nagoya Math.\ J.\ 237 (2020),
127--165, Section~1.2, equation~(6) and Theorem~1, printed pp.~133--134,
does treat a series whose summands contain
\[
 L_{-d}\bigl(\sigma-(k+1)/2\bigr)
 L_{-d}\bigl(\eta-(k+1)/2\bigr).
 \tag{HM.2}
\]
Here the $L_D$ notation includes the nonfundamental-discriminant divisor
factors.  The theorem gives meromorphic continuation in three complex
variables, an explicit holomorphic regularization, and a vector functional
equation.  The parameters exclude poles of the summands before continuation.
Proposition~3 and its complete proof, printed pp.~152--153, and the
complete proof of Theorem~1, pp.~153--154, establish these assertions by
regularized integrals and an invertible matrix of Whittaker integrals.
They do not identify the coefficients in (HM.2) with a holomorphic
modular $q$-series over a specified $p$-integral lattice.  In particular,
this cited theorem contains neither an integral residue formula forcing
one coefficient to be a unit, nor a bound $d<p$.  A nonzero complex
residue alone supplies neither conclusion.  The partial Euler product
in Section~3.4, p.~154, sums the square-discriminant subseries; it is not
an Euler product for the negative fundamental discriminants in (HM.1).

Mizuno's separate \emph{Petersson norms of Eisenstein series and
Kohnen--Zagier's formula}, J.\ Th\'eor.\ Nombres Bordeaux 32 (2020), 665--684,
Theorem~2.4, p.~670, gives the regularized Kohnen--Zagier formula for
\emph{one} Cohen Eisenstein coefficient squared.  For a fundamental $D$
with $(-1)^rD>0$, its proof in Section~5, pp.~680--681, uses
\[
 L(G_{2r},D,s)=L(s,\chi_D)L(s+1-2r,\chi_D)
 \tag{HM.3}
\]
and the functional equation to obtain the square of $L(1-r,\chi_D)$.
The complete proof and the norm-ratio normalization have been read.
This is not the mixed product of the weights $b$ and $3$ in (HM.1).
No $p$-unit conclusion follows just from the nonzero real regularized
norms.  The earlier computations of those norms are named inputs to
that proof, not newly rechecked in this audit.

\paragraph{A saturated Rankin--Cohen trial still fails the diagonal test.}
The standard Rankin--Cohen formula is a bilinear differential expression,
so its Fourier coefficients are additive convolutions.  One might try
order $p-1$ to select equal indices by Fermat's theorem, while keeping
the formal weight linear in $p$.  The following exact calculation rules
out this particular trial, even after its common factor $p$ is removed.
It is a statement about the displayed universal operator; it does not
exclude an accidental identity for the particular Cohen series.

\begin{proposition}[KT-SATURATED-RANKIN-COHEN-SCOPE]
Set $k=(p+4)/2$, $\ell=7/2$, and $N=p-1$, and define in
$\mathbb Z_{(p)}$
\[
 c_r=(-1)^r\binom{N+k-1}{N-r}\binom{N+\ell-1}{r}
     =(-1)^r\binom{3p/2}{p-1-r}\binom{p+3/2}{r},
 \qquad 0\le r\le p-1.
 \tag{HM.4}
\]
Every $c_r$ is divisible by $p$.  For $s_r=c_r/p$ one has
\[
 s_r\equiv0\pmod p\quad(b+1\le r\le p-2),
 \qquad
 s_r\not\equiv0\pmod p\quad(0\le r\le b\text{ or }r=p-1).
 \tag{HM.5}
\]
Put $S(X,Y)=\sum_{r=0}^{p-1}\bar s_rX^rY^{p-1-r}$.
There are no scalars $\alpha,\beta\in\mathbb F_p$ such that, for all
$x,y\in\mathbb F_p$ with $x+y\ne0$,
\[
 \alpha S(x,y)+\beta=\begin{cases}1,&x=y,\\0,&x\ne y.\end{cases}
 \tag{HM.6}
\]
Thus this saturated operator, even with an arbitrary scalar multiple
of the ordinary product added, is not a universal equal-index selector.
\end{proposition}

\begin{proof}
All factorial denominators in (HM.4) have indices below $p$.
If $s=p-1-r>0$, the numerator of $\binom{3p/2}{s}$, after multiplying
by the unit $2^s$, is
$\prod_{j=0}^{s-1}(3p-2j)$.  Exactly its $j=0$ factor is divisible by
$p$, and it has valuation one.  For $s=0$ the binomial is $1$.
The numerator of the other binomial, after multiplying by $2^r$, is
$\prod_{j=0}^{r-1}(2p+3-2j)$.  It contains a multiple of $p$ precisely
when $r\ge b+1$, and that multiple is its $j=b$ factor, equal to $p$.
Its valuation is then one.  This proves the exact valuations and (HM.5).

For $x+y\ne0$ the desired indicator is the value of the homogeneous
polynomial
\[
 D(X,Y)=(X+Y)^{p-1}-(X-Y)^{p-1}
       =\sum_{r=0}^{p-1}\bigl((-1)^r-1\bigr)X^rY^{p-1-r}
       \quad\text{over }\mathbb F_p.
 \tag{HM.7}
\]
If (HM.6) held, the homogeneous polynomial
$\alpha S+\beta(X+Y)^{p-1}-D$ would vanish on every point of the line
$X+Y=1$.  Its restriction to that line has degree at most $p-1$ and
has $p$ roots, so is the zero polynomial.  Homogenization recovers the
original homogeneous polynomial, which is therefore zero.
The coefficients of $X^{p-3}Y^2$ and $X^{p-2}Y$ now contradict each
other: both corresponding coefficients of $S$ vanish by (HM.5), while
those of $(X+Y)^{p-1}$ are $1,-1$ and those of $D$ are $0,-2$.
They would force $\beta=0$ and $\beta=2$.
\end{proof}

Writing $\Theta=q\,d/dq$, the formal operator just tested is
\[
 \mathcal R(f,g)=\sum_{r=0}^{p-1}s_r
          \Theta^rf\,\Theta^{p-1-r}g.
 \tag{HM.8}
\]
Its symbol has been checked on arbitrary residue indices.  The proposition
does not assert that a pair of genuine modular forms can have arbitrary
Fourier coefficients, and does not prove that no other bilinear operator
or global modular relation works.  If a compatible modular realization
is used, adding a multiple of $A^2fg$, with $A\equiv1\pmod p$ a Hasse
form of weight $p-1$, has exactly the additional symbol $\beta$ already
excluded.  Restricting to the actual Cohen forms or to a finite short
window could impose extra relations not addressed by this calculation.

The source formula for the trial is Zagier, \emph{Modular forms and
differential operators}, Proc.\ Indian Acad.\ Sci.\ (Math.\ Sci.)
104 (1994), 57--75, Section~1, equation~(1), p.~58.  The entire Section~1
and both modularity explanations on pp.~58--59 were read; the displayed
binomial formula was also visually checked.  That paper's opening
convention is integral-weight modular forms.  The proposition above is
pure coefficient algebra with the specified half-integral parameters,
and does not silently import a half-integral-weight theorem from that
convention.  No modularity assertion is needed for its obstruction.

\paragraph{Finite square-class support is a different hypothesis.}
Ahlgren--Choi--Rouse, \emph{Congruences for level four cusp forms},
Math.\ Res.\ Lett.\ 16 (2009), author Theorem~1.2, pp.~2--3, assumes a
nonzero cusp form in the level-four plus space, integral at a place
above a prime $p\ge5$, whose reduction is supported on \emph{finitely
many} square classes.  Under
$\lambda+1/2<p(p+1+1/2)$ it concludes that $\lambda$ is even and gives
a derivative-of-theta formula.  The complete stated theorem and its
final proof on author pp.~16--17 were read; the intermediate
propositions used in that proof are not newly certified here.
The weight bound is not the decisive missing premise in (HM.1).
$H_b$ is Eisenstein, and the assertion that its support lies in
$\{D:\beta_3(\chi_D)=0\bmod p\}$ does not imply that this set is finite.
Passing to a derivative would not establish that missing finiteness.
This theorem therefore does not convert the hypothetical disjoint
supports into a contradiction.

\paragraph{Evidence and boundary.}
The four primary PDFs and their exact retrieval hashes are saved in
\texttt{evidence/lean\_chain/primary/hadamard\_scope/}.
The separate standard-library script
\texttt{saturated\_rankin\_cohen\_control.py} checks all coefficients
in (HM.4)--(HM.5) for the single diagnostic prime $31$, including the
two conflicting coefficients in (HM.7).  This is a finite algebra
control, not a prime census or a proof by computation of (HM.1).
No simultaneous seed, original class-component vanishing, or new
finite $K$-group degree is inferred in this section.
